% 5 -------------------------------------------------------------------------
\begin{frame}{A molecular energy is only the starting point}
\centering
\begin{tikzpicture}
\node[card,text width=2.45cm,minimum height=1.5cm](a) at (0,0){\textbf{Electronic}\\Bonding and\\3D geometry};
\node[card,text width=2.45cm,minimum height=1.5cm](b) at (3.45,0){\textbf{Thermal}\\Vibrations and\\entropy};
\node[card,text width=2.45cm,minimum height=1.5cm](c) at (6.9,0){\textbf{Solvent}\\Stabilization\\in water};
\node[card,text width=2.45cm,minimum height=1.5cm](d) at (10.35,0){\textbf{Chemical states}\\Protonation and\\hydration};
\node[font=\Large] at (1.72,0){+};\node[font=\Large] at (5.17,0){+};\node[font=\Large] at (8.62,0){+};
\node[card,text width=5.5cm,minimum height=.85cm](e) at (5.17,-2.25){\textbf{Biochemical reaction Gibbs energy}};
\draw[arr](a.south)--(e.north west);\draw[arr](b.south)--(e.north);\draw[arr](c.south)--(e.north);\draw[arr](d.south)--(e.north east);
\end{tikzpicture}
\takeaway{The pipeline connects microscopic chemistry to a biochemical quantity.}
\slidenote{UMA is a machine-learned interatomic potential: it predicts energies and forces from 3D structures.}
\end{frame}

% 6 -------------------------------------------------------------------------
\begin{frame}{What quantity does MetaG predict?}
\centering
{\LARGE $\displaystyle\Delta_rG^\prime=\underbrace{\dg}_{\text{MetaG output}}+\underbrace{RT\ln Q^\prime}_{\text{concentration effect}}$}\par\vspace{.45cm}
\begin{columns}[T]
\column{.48\textwidth}
\begin{block}{Standard biochemical conditions}
\small 298.15 K \quad $\vert$ \quad pH 7\\[.2cm]
$I=0$ \quad $\vert$ \quad no Mg binding
\end{block}
\column{.48\textwidth}
\begin{block}{Reaction as written}
\small 1 M solutes; pure liquid water\\[.2cm]
Energy reported in kJ/mol
\end{block}
\end{columns}
\gap
\takeaway{Use the standard energy with concentrations to assess a pathway.}
\slidenote{These conditions are fixed in the pipeline. $Q^\prime$ is the biochemical reaction quotient.}
\end{frame}

% 7 -------------------------------------------------------------------------
\begin{frame}{The pipeline: structures in, reaction energy out}
\centering
\begin{tikzpicture}[x=1cm,y=1cm]
\node[card,text width=3.5cm,minimum height=1.2cm](a) at (0,0){\textbf{1. Validate}\\Structures and balance};
\node[card,text width=3.5cm,minimum height=1.2cm](b) at (4.6,0){\textbf{2. Represent}\\Full molecules or cores};
\node[card,text width=3.5cm,minimum height=1.2cm](c) at (9.2,0){\textbf{3. Sample}\\3D conformer minima};
\node[card,text width=3.5cm,minimum height=1.2cm](d) at (9.2,-2.1){\textbf{4. Compute}\\Species free energies};
\node[card,text width=3.5cm,minimum height=1.2cm](e) at (4.6,-2.1){\textbf{5. Assemble}\\Stoichiometry and pH};
\node[card,text width=3.5cm,minimum height=1.2cm](f) at (0,-2.1){\textbf{6. Report}\\$\dg$ and diagnostics};
\draw[arr](a)--(b);\draw[arr](b)--(c);\draw[arr](c)--(d);\draw[arr](d)--(e);\draw[arr](e)--(f);
\end{tikzpicture}
\gap\takeaway{Each prediction retains the molecular representation and energy terms.}
\end{frame}

% 8 -------------------------------------------------------------------------
\begin{frame}{One molecule can occupy several 3D conformations}
\centering
\begin{tikzpicture}
\node[card,text width=3.25cm,minimum height=1.4cm](a) at (0,0){Generate a\\conformer pool};
\node[card,text width=3.25cm,minimum height=1.4cm](b) at (4.4,0){Rank and relax\\with UMA};
\node[card,text width=3.25cm,minimum height=1.4cm](c) at (8.8,0){Validate minima\\and combine};
\draw[arr](a)--(b);\draw[arr](b)--(c);
\node[text width=3.4cm,align=center,font=\small] at (0,-1.65){Different starting shapes};
\node[text width=3.4cm,align=center,font=\small] at (4.4,-1.65){Energies and forces};
\node[text width=3.4cm,align=center,font=\small] at (8.8,-1.65){Thermal and solvent terms};
\end{tikzpicture}
\gap
\takeaway{An ensemble includes both low-energy structures and their populations.}
\slidenote{Sampling expands until the calculated ensemble stabilizes. A plateau does not bound undiscovered conformers.}
\end{frame}

% 9 -------------------------------------------------------------------------
\begin{frame}{From conformers to species, then to reactions}
\begin{columns}[c]
\column{.48\textwidth}
\centering\textbf{Within a species}\gap
$g_k=E_k^{\rm UMA}+G_k^{\rm thermal}+\Delta G_k^{\rm solvent}$\gap
$\displaystyle G_i=-RT\ln\sum_k e^{-g_k/RT}+C^\circ$\gap
\small Low-energy conformers carry more weight.
\column{.48\textwidth}
\centering\textbf{Across a reaction}\gap
$\displaystyle\dg=\sum_i\nu_iG_i^{\rm eff}+C_{\rm proton/pH}$\gap
\small Products add; reactants subtract.\\[.15cm]
$G_i^{\rm eff}$ also includes enabled hydrate states.
\end{columns}
\gap\takeaway{Every reaction energy can be traced back to its molecular contributions.}
\slidenote{$C^\circ$: standard-state correction. The appendix gives the proton, water, and pKa conventions.}
\end{frame}

% 10 ------------------------------------------------------------------------
\begin{frame}{Large cofactors: retain the chemistry that changes}
\centering\includegraphics[width=.97\linewidth,height=4.25cm,keepaspectratio]{\CoreFigure}
\takeaway{The structural core preserves the reactive thioester and its local scaffold.}
\slidenote{CoA example from the router. Balance and removed-fragment identity are checked; cores remain approximations.}
\end{frame}

% 11 ------------------------------------------------------------------------
\begin{frame}{The aqueous species is a mixture of chemical states}
\centering
\begin{tikzpicture}
\node[card,minimum width=3.1cm,minimum height=.85cm](a) at (0,0){Acid form};
\node[card,minimum width=3.1cm,minimum height=.85cm](b) at (6.4,0){Ionized form};
\node[card,minimum width=3.1cm,minimum height=.85cm](c) at (0,-1.65){Carbonyl form};
\node[card,minimum width=3.1cm,minimum height=.85cm](d) at (6.4,-1.65){Hydrated form};
\draw[<->,teal,thick](a)--node[above]{pH and pKa}(b);
\draw[<->,teal,thick](c)--node[above]{Water and $K_{\rm hyd}$}(d);
\end{tikzpicture}
\gap
\takeaway{State populations shift the effective free energy of a metabolite.}
\slidenote{Production uses independent pKa sites and enabled carbonyl hydration. The two equilibria are schematic; coupled-state routing is not assumed.}
\end{frame}

\begin{frame}{Current scoring and the new validation tools}
\begin{columns}[T]
\column{.48\textwidth}
\begin{block}{Active reaction scorer}
\small UMA + gas thermal terms + ALPB.\gap
Tabulated pKa patterns; largely independent sites.\gap
Unsupported isotope inputs are withheld; the legacy SMD reaction correction is blocked.
\end{block}
\column{.48\textwidth}
\begin{block}{Implemented research tools}
\small Coupled-state registry: references, conditions, and thermodynamic cycles.\gap
Aqueous-ensemble admission and network/covariance audits.\gap
These tools do not replace the default energy or pKa model.
\end{block}
\end{columns}
\takeaway{The package now checks more assumptions; its default pKa physics is unchanged.}
\slidenote{Integration snapshot: 8 October 2026. Registry transitions require justified state-specific data; the registry does not predict missing pKa values.}
\end{frame}
